\chapter{Prime-Fiber Zoom Renormalization Tower}
\label{ch:visions-rh-prime-fiber-zoom-renormalization-tower}
\origin{ai}

\autoref{ch:visions-rh-memory-filtration-prime-zeckendorf-register}
separates append-only history from orthogonal residual readout, and
\autoref{ch:visions-rh-pzg-phase-readback-log-energy} assigns PZG
coordinates a logarithmic energy. The present chapter records the
renormalization surface that connects these two rows. A phase-relevant
observation selects a prime-fiber cell, renormalizes that cell to the
current visible surface, and records every unselected component in a gap
ledger.

\section{Prime-Fiber Zoom Systems}
\label{sec:rh-prime-fiber-zoom-systems}

\begin{definition}[Prime-fiber zoom system]
\label{def:rh-prime-fiber-zoom-system}
A prime-fiber zoom system consists of a base fibre $\mathcal F_e$, finite
address sets
$$
A_m\subseteq \PrimeUp\times\mathbb N,
$$
cell families
$$
\mathcal P_m:=\{C_{\alpha}:\alpha\in A_m\},
$$
selection rows
$$
\operatorname{Sel}_{\alpha}:\mathcal F_e\to C_{\alpha},
$$
zoom rows
$$
\operatorname{Zoom}_{\alpha}:C_{\alpha}\to\mathcal F_e,
$$
and an orthogonal register basis $(e_{\alpha})_{\alpha}$ such that
different displayed addresses have orthogonal basis rows. A step with
address $\alpha_m$ is the composite
$$
T_{\alpha_m}
:=
\operatorname{Zoom}_{\alpha_m}\circ\operatorname{Sel}_{\alpha_m}.
$$
The unselected part of $\mathcal F_e$ is not discarded; it is exported
to the gap row governed by
\autoref{prin:ask-sig-spine-compression-memory}.
\end{definition}

\begin{definition}[Solenoid clopen prime cell]
\label{def:rh-solenoid-clopen-prime-cell}
Fix a finite displayed prime family $S$ and the solenoid phase channel
$$
\pi_S:\Sigma_S\to\mathbb T
$$
of \autoref{def:rh-phase-channel-hidden-prime-fibre}. Its legal hidden
fibre is
$$
K_S:=\ker(\pi_S).
$$
A solenoid clopen prime cell is a finite fibre-coordinate row
$$
C(\mathbf a,\mathbf k)
:=
\prod_{p\in S}\bigl(a_p+p^{k_p}\mathbb Z_p\bigr)
$$
inside a displayed product presentation of the finite prime-register
window, where $k_p\ge0$ and $a_p$ is read modulo $p^{k_p}$. The cell
records a finite prefix of the hidden prime-register state over the
visible phase point.
\end{definition}

\begin{definition}[Prime-fibre zoom row]
\label{def:rh-prime-fibre-zoom-row}
For a solenoid clopen prime cell
$C(\mathbf a,\mathbf k)$, the prime-fibre zoom row is the componentwise
renormalization
$$
\operatorname{Zoom}_{\mathbf a,\mathbf k}
:
C(\mathbf a,\mathbf k)\to
\prod_{p\in S}\mathbb Z_p
$$
given by
$$
\operatorname{Zoom}_{\mathbf a,\mathbf k}
\bigl((x_p)_{p\in S}\bigr)
:=
\left(
\frac{x_p-a_p}{p^{k_p}}
\right)_{p\in S}.
$$
Its inverse inclusion is
$$
\operatorname{Cell}_{\mathbf a,\mathbf k}
\bigl((y_p)_{p\in S}\bigr)
:=
\bigl(a_p+p^{k_p}y_p\bigr)_{p\in S}.
$$
Thus a selected fibre cell is read as a whole prime-register window
after renormalization.
\end{definition}

\begin{theorem}[Solenoid clopen refinement gives an orthogonal filtration]
\label{thm:rh-solenoid-clopen-refinement-orthogonal-filtration}
Let
$$
C_0\supseteq C_1\supseteq\cdots\supseteq C_N
$$
be a finite nested sequence of solenoid clopen prime cells over the same
finite prime family $S$. Let $\mathcal A_m$ be the finite sigma-algebra
generated by the stage-$m$ cell partition of $K_S$, and set
$$
\mathcal H_m:=L^2(K_S,\mathcal A_m).
$$
Then
$$
\mathcal H_m\subseteq\mathcal H_{m+1},
\qquad
\mathcal H_{m+1}
=
\mathcal H_m\oplus
\bigl(\mathcal H_{m+1}\ominus\mathcal H_m\bigr).
$$
Consequently, for
$$
\eta_{m+1}\in
\mathcal H_{m+1}\ominus\mathcal H_m
$$
and $v_m\in\mathcal H_m$,
$$
\|v_m+\eta_{m+1}\|^2
=
\|v_m\|^2+\|\eta_{m+1}\|^2.
$$
\end{theorem}

\begin{proof}
Refinement of finite cell partitions gives
$\mathcal A_m\subseteq\mathcal A_{m+1}$, hence
$\mathcal H_m\subseteq\mathcal H_{m+1}$. The complement
$\mathcal H_{m+1}\ominus\mathcal H_m$ is defined as the orthogonal
complement inside the Hilbert row $\mathcal H_{m+1}$. The displayed
direct sum and norm identity are therefore the Hilbert projection
identity, specialized to the solenoid clopen-cell filtration.
\end{proof}

\begin{definition}[Zoom history tower]
\label{def:rh-prime-fiber-zoom-history-tower}
A zoom history tower is a finite or countable sequence of addresses
$$
\alpha_0,\alpha_1,\ldots
\qquad
(\alpha_j\in\PrimeUp\times\mathbb N)
$$
with nested cells
$$
C_{\alpha_0}
\supseteq
C_{\alpha_0\alpha_1}
\supseteq
C_{\alpha_0\alpha_1\alpha_2}
\supseteq\cdots.
$$
The history row is append-only:
$$
h_{m+1}=h_m\alpha_m.
$$
The visible row at stage $m+1$ is the zoomed cell
$$
B_{m+1}
:=
\operatorname{Zoom}_{\alpha_m}(B_m\cap C_{\alpha_m}),
$$
while the Hilbert residual row is the address vector $e_{\alpha_m}$.
\end{definition}

\begin{theorem}[Prime-fiber zoom gives history, similarity, and orthogonality]
\label{thm:rh-prime-fiber-zoom-history-similarity-orthogonality}
In a prime-fiber zoom history tower, each displayed step has the three
rows
$$
h_m\subseteq h_{m+1},
$$
$$
B_{m+1}
=
\operatorname{Zoom}_{\alpha_m}(B_m\cap C_{\alpha_m}),
$$
and
$$
e_{\alpha_m}\perp
\operatorname{span}\{e_{\alpha_0},\ldots,e_{\alpha_{m-1}}\}.
$$
Consequently, if
$$
R_m:=\sum_{j<m}e_{\alpha_j},
\qquad
R_{m+1}:=R_m+e_{\alpha_m},
$$
then
$$
\|R_{m+1}\|^2
=
\|R_m\|^2+\|e_{\alpha_m}\|^2.
$$
\end{theorem}

\begin{proof}
The history inclusion is the append-only clause of
\autoref{def:rh-prime-fiber-zoom-history-tower}, read through the same
history surface as
\autoref{def:kernel-spine-history-carrier} and
\autoref{def:kernel-spine-continuation}. The displayed formula for
$B_{m+1}$ is the definition of the zoomed visible row. Orthogonality is
part of the register basis in
\autoref{def:rh-prime-fiber-zoom-system}; hence the final identity is
the Pythagorean identity for the decomposition
$R_{m+1}=R_m+e_{\alpha_m}$.
\end{proof}

\begin{principle}[Prime-fiber zoom is not erasure]
\label{prin:rh-prime-fiber-zoom-not-erasure}
The zoom operation changes the current visible surface, but it does not
license unledgered deletion. The selected cell is renormalized to the
next visible row, the address $\alpha_m$ is appended to the history row,
and the outside of the selected cell is represented by the gap ledger.
Thus
$$
B_{m+1}\supseteq B_m,
\qquad
B_{m+1}\sim_{\exp} B_m,
\qquad
B_{m+1}/B_m\perp B_m
$$
are not competing assertions. They are the history, zoom, and residual
projections of one certified step.
\end{principle}

\section{PZG Energy Compatibility}
\label{sec:rh-prime-fiber-zoom-pzg-energy}

\begin{definition}[PZG zoom step]
\label{def:rh-pzg-zoom-step}
A PZG zoom step is a prime-fiber zoom step whose address is
$$
\alpha=(p,k)
$$
and whose residual vector is the log-energy basis vector
$\varepsilon_{p,k}$ of
\autoref{def:rh-log-energy-pzg-readout}. Its squared residual energy is
$$
\|\varepsilon_{p,k}\|^2=F_k\log p.
$$
\end{definition}

\begin{principle}[PZG zoom carries preserve triangular energy]
\label{prin:rh-pzg-zoom-carry-preserves-triangular-energy}
If two consecutive PZG zoom addresses in one prime row create a
Zeckendorf adjacency, the legal normal form replaces
$$
\varepsilon_{p,k}+\varepsilon_{p,k+1}
\rightsquigarrow
\varepsilon_{p,k+2}.
$$
By \autoref{prin:rh-zeckendorf-carries-preserve-log-energy}, this
replacement preserves log-energy. A PZG zoom step is therefore not an
arbitrary bit change in a hidden fibre; it is a prime-indexed residual
update compatible with the triangular square-norm row.
\end{principle}

\section{Analytic Prime-Shift Constraint}
\label{sec:rh-analytic-prime-shift-constraint}

\begin{definition}[Analytic prime-shift tower]
\label{def:rh-analytic-prime-shift-tower}
An analytic prime-shift tower is a prime-fiber zoom history tower
together with finite holomorphic readouts
$$
F_m:\Omega_m\to\ComplexUp
$$
and ledgers $\mathcal L_m$ such that the address sequence
$$
\alpha_0,\alpha_1,\ldots
\qquad
(\alpha_m=(p_m,k_m))
$$
is not merely a discrete stream of prime coordinates. It must satisfy
the following certificate rows:
$$
\begin{gathered}
\mathsf{NestedCell},\quad
\mathsf{OrthogonalOrZeckendorfCarry},\quad
\mathsf{HeightControl},\\
\mathsf{HolomorphicChartCompatibility},\quad
\mathsf{DominantStrandCompatibility},\quad
\mathsf{RoucheNoWallCrossing},\quad
\mathsf{MonodromyLedger}.
\end{gathered}
$$
The nested-cell row says that the selected cells form a refinement chain.
The orthogonal-or-carry row says that each residual address is either
orthogonal to the previous Hilbert row or is absorbed by a legal
Zeckendorf carry. The height-control row supplies compact-uniform
control for the exponential reads
$$
\exp(-sH_m),
\qquad
H_m:=\sum_{j\le m}\|e_{\alpha_j}\|^2,
$$
or a displayed regularization producing the same compact-uniform
holomorphic limit. The holomorphic-chart row supplies analytic
continuation data in the sense of
\autoref{thm:analytic-continuation-operation}; the dominant-strand row
records the branch selected by the tower; the Rouch\'e row records
zero-count stability on displayed boundaries; and the monodromy ledger
records branch return data instead of treating it as invisible motion.
\end{definition}

\begin{definition}[SmoothPrimeShiftConstraint]
\label{def:rh-smooth-prime-shift-constraint}
The obligation
$$
\mathsf{SmoothPrimeShiftConstraint}
$$
asserts that every zero-relevant prime-fiber index stream used by the
completed zeta route is an analytic prime-shift tower in the sense of
\autoref{def:rh-analytic-prime-shift-tower}. It therefore rules out an
unconstrained sequence of addresses
$$
\alpha_0,\alpha_1,\alpha_2,\ldots
$$
as a source for a holomorphic zeta observation. A legal stream must
carry nesting, orthogonality or carry, height control, chart
compatibility, branch ledger, no-wall-crossing, and cofinal coverage.
\end{definition}

\begin{theorem}[Analytic prime-shift towers give holomorphic readouts]
\label{thm:rh-analytic-prime-shift-tower-holomorphic-readout}
Let an analytic prime-shift tower have domains
$$
\Omega_m\subseteq\Omega_{m+1},
\qquad
\Omega:=\bigcup_m\Omega_m,
$$
with holomorphic readouts $F_m$ on $\Omega_m$. Assume that the chart
compatibility rows identify the restrictions on overlaps and that
$F_m\to F$ uniformly on every compact subset of $\Omega$. Then $F$ is
holomorphic on $\Omega$, with branch and zero-island data controlled by
the monodromy and Rouch\'e ledgers of the tower.
\end{theorem}

\begin{proof}
On each compact subset $K\subseteq\Omega$, the compact-uniform convergence
row gives uniform convergence of holomorphic functions on a neighbourhood
of $K$. The Weierstrass theorem gives a holomorphic limit on that
neighbourhood. The holomorphic-chart compatibility rows identify these
local limits on overlaps by the analytic-continuation operation. The
monodromy ledger records possible branch return data, so no branch
selector is silently changed. The Rouch\'e no-wall-crossing row records
the boundary inequalities needed to preserve zero counts inside the
displayed zero islands. These rows together give the claimed holomorphic
readout on $\Omega$.
\end{proof}

\begin{principle}[Prime shifts that are not analytic remain address streams]
\label{prin:rh-nonanalytic-prime-shifts-remain-address-streams}
A PZG address stream that lacks the rows of
\autoref{def:rh-analytic-prime-shift-tower} may still encode discrete
prime-register information. It does not thereby generate a continuous
curve, a differentiable observation, or a holomorphic zeta readout.
Analytic operations such as eta continuation, functional-equation
transport, contour transport, or explicit-formula transport are precisely
the rows that convert a controlled prime-register stream into an
analytic observation object.
\end{principle}

\section{Continuous Exponential Projection}
\label{sec:rh-prime-fiber-zoom-continuous-exponential-projection}

\begin{definition}[Internal height and external exponential readout]
\label{def:rh-internal-height-external-exponential-readout}
The internal height of a PZG register is the additive square-norm row
$$
H(R):=\|R\|^2.
$$
For the integer register $R_{\mathsf{PZG}}(n)$ of
\autoref{def:rh-log-energy-pzg-readout}, this is
$$
H(R_{\mathsf{PZG}}(n))=\log n
$$
by \autoref{thm:rh-log-energy-reads-logarithmic-height}. The external
exponential readout is the multiplicative row
$$
\operatorname{Scale}(R):=\exp(H(R)).
$$
The endpoint symbol $\mathbf e$ denotes the identity endpoint or base
fibre in the zoom system; the scalar $\exp(1)$ is the unit value of the
continuous exponential readout.
\end{definition}

\begin{definition}[Character-valued exponential readout]
\label{def:rh-character-valued-exponential-readout}
Let $A$ be a displayed additive ledger group, and let
$$
\omega:A\to(\mathbb C,+)
$$
be an additive readout weight. The exponential readout determined by
$\omega$ is
$$
\operatorname{Exp}_{\omega}:A\to\mathbb C^\times,
\qquad
\operatorname{Exp}_{\omega}(x):=\exp(\omega(x)).
$$
It is the external multiplicative observation of the internal additive
coordinate. If $\omega(x)\in i\RealUp$ for every displayed $x$, then
$$
\operatorname{Exp}_{\omega}(x)\in\mathbb T
$$
and the readout is a unit-circle phase projection. If
$\omega(x)\in\RealUp$ for every displayed $x$, then
$$
\operatorname{Exp}_{\omega}(x)\in\RealUp_{>0}
$$
and the readout is a radial-scale projection. Thus the same exponential
mechanism separates into the two rows
$$
\exp(a+ib)=\exp(a)\exp(ib),
\qquad
\exp(ib)\in\mathbb T,
\qquad
\exp(a)\in\RealUp_{>0}.
$$
\end{definition}

\begin{theorem}[Exponential readout is the normalized additive-to-multiplicative projection]
\label{thm:rh-exponential-readout-normalized-additive-multiplicative}
Let
$$
\Phi:\RealUp\to\RealUp_{>0}
$$
be a differentiable readout satisfying
$$
\Phi(t+s)=\Phi(t)\Phi(s),
\qquad
\Phi(0)=1,
\qquad
\Phi'(0)=1.
$$
Then
$$
\Phi(t)=\exp(t).
$$
\end{theorem}

\begin{proof}
Set
$$
f(t):=\log\Phi(t).
$$
The multiplicative law for $\Phi$ gives
$$
f(t+s)=f(t)+f(s).
$$
Differentiability gives $f(t)=ct$ for a constant $c$. Since
$$
\Phi'(0)=f'(0)\exp(f(0))=c,
$$
the normalization $\Phi'(0)=1$ gives $c=1$. Hence
$\Phi(t)=\exp(t)$.
\end{proof}

\begin{principle}[Pythagorean addition projects to exponential multiplication]
\label{prin:rh-pythagorean-addition-exponential-multiplication}
If a PZG zoom step has
$$
R_{m+1}=R_m+e_{\alpha_m},
\qquad
e_{\alpha_m}\perp R_m,
$$
then
$$
H(R_{m+1})=H(R_m)+\|e_{\alpha_m}\|^2.
$$
Under the external exponential readout,
$$
\operatorname{Scale}(R_{m+1})
=
\operatorname{Scale}(R_m)\exp(\|e_{\alpha_m}\|^2).
$$
Thus the internal triangular square-norm row is additive, while its
continuous scale projection is multiplicative.
\end{principle}

\begin{principle}[Solenoid exponential readout is phase-valued]
\label{prin:rh-solenoid-exponential-readout-phase-valued}
For the localized additive ledger
$$
A=\mathbb Z[S^{-1}],
$$
the visible solenoid phase coordinate is read through unit characters of
the form
$$
q\mapsto\exp(2\pi i q\theta)\in\mathbb T.
$$
This is the instance of
\autoref{def:rh-character-valued-exponential-readout} with
$$
\omega_{\theta}(q)=2\pi i q\theta.
$$
The hidden row of
\autoref{def:unified-reality-s-solenoid-over-unit-circle} is the compact
prime-register fibre $K_S$, not a positive real scale group. Therefore a
real exponential row
$$
q\mapsto\exp(aq),
\qquad
a\in\RealUp,
$$
is radial scale data rather than solenoid-hidden phase data, unless
$a=0$ on the zero-relevant packet.
\end{principle}

\begin{principle}[Phase and radial reads use the same internal height]
\label{prin:rh-phase-radial-same-internal-height}
For a zeta monomial and $s=\frac12+r+it$,
$$
n^{1/2-s}
=
\exp\bigl(-(r+it)\log n\bigr)
=
\exp(-rH(R_{\mathsf{PZG}}(n)))
\exp(-itH(R_{\mathsf{PZG}}(n))).
$$
Equivalently, in PZG coordinates,
$$
n^{1/2-s}
=
\exp\left(
-(r+it)\sum_{p,k}z_{p,k}F_k\log p
\right).
$$
The second factor is the unit-circle phase read; the first factor is the
positive radial-scale read. Both consume the same internal PZG height.
Only the unit-circle factor is a phase projection. A nontrivial radial
factor is an external scale projection and must be ledgered as such.
\end{principle}

\begin{principle}[RH reading of prime-fiber zoom]
\label{prin:rh-prime-fiber-zoom-rh-reading}
For a zeta phase row
$$
n^{-it}
=
\exp\left(
-it\sum_{p,k}z_{p,k}F_k\log p
\right),
$$
the PZG coordinates are prime-fiber zoom addresses read through phase.
For a centered zero packet with
$$
z_{\rho}=1-\frac1\rho=e^{\ell+i\theta},
$$
the angular coordinate $\theta$ is compatible with unit phase readout,
whereas a nonzero radial coordinate $\ell$ would have to appear as a
ledgered residual, a gap row, or a refused scale coordinate. It cannot
be absorbed into the visible wave without a prime-fiber address or a
legal solenoid fibre row.
\end{principle}
